\scanpage{68}
% archive sheet number: 69
conserve ces structures préhilbertiennes \deleted{et} Noter que à une décomposition de $E$ \deleted{\unsure{en somme directe}} en somme de \unsure{sous-espaces} mutuellement orthogonaux $E_0$ et $E_1$ donne une décomposition orthogonale

\[
\overset{n}{\Lambda}^{(2)}E=\sum_{0\leq p\leq n}\overset{n-p}{\Lambda}^{(2)}E_0\overset{(2)}{\otimes}\overset{p}{\Lambda}^{(2)}E_1 .
\]

Si $u$ est un op. dans $E$ s'annulant sur $E_1$, alors on a $(\overset{n}{\Lambda}u).(x_1\wedge\ldots\wedge x_n)\inserted{=ux_1\wedge\ldots\wedge ux_n}=0$, \deleted{\unsure{pourvu|puisque}} dès que l'un \deleted{\unsure{des}} des $x_i$ \deleted{\unsure{dans} $\overset{n}{\Lambda}$} $\unsure{\in} E_1$, donc $\overset{n}{\Lambda}u$ s'annule sur \deleted{\unsure{le}} toutes les \deleted{\unsure{composantes}} \inserted{termes composants} de $\sum_p\overset{n-p}{\Lambda}E_0\otimes\overset{p}{\Lambda}E_1$ sauf le premier, \deleted{et si de même \unsure{soit}}

il en résulte par prolongement par continuité que l'op. correspondant dans $\overset{n}{\Lambda}^{(2)}E$ s'annule sur tous les $\overset{n-p}{\Lambda}^{(2)}E_0\overset{(2)}{\otimes}\overset{p}{\Lambda}^{(2)}E_{\unsure{1|i}}$, \deleted{\unsure{sauf}} avec $p\neq0$. \deleted{D\unsure{e}} \inserted{On voit de} même que, si $u$ invarie $E_0$, alors $\overset{n}{\Lambda}u$ conserve $\overset{n}{\Lambda}^{(2)}E_0$, de plus dans \deleted{\unsure{cet}} \deleted{\unsure{esp|cet}} ce dernier espace, il se réduit à

$\overset{n}{\Lambda}v$, où $v$ est la restriction de $u$ à $E_0$. Dans le cas actuel, \deleted{\illegible{2}} \inserted{($u(E_1)=0$, $u(E_0)\subset E_0$, \deleted{de} $E_0$ de dim. finie)} on voit donc que \deleted{$\overset{n}{\Lambda}u$ \unsure{se ramène}} les valeurs propres de $\overset{n}{\Lambda}u$ dans $\overset{n}{\Lambda}^{(2)}E$, sont les valeurs propres de $\overset{n}{\Lambda}v$ dans $\uline{\overset{n}{\Lambda}^{(2)}E_0}=\overset{n}{\Lambda}E_0$. Il est bien connu, que

\deleted{\unsure{si est de}} \deleted{ce} puisqu'il s'agit d'opérateurs dans l'espace vectoriel de dim. finie $E_0$, que ces valeurs propres s'obtiennent en \deleted{\unsure{prenant}} \inserted{tous} les produits \inserted{de $n$} \deleted{\unsure{au}} valeurs propres de $v$ correspondant à des indices distincts (on le voit par ex en mettant $v$ sous forme de matrice triangulaire\unsure{?|/}\unsure{ainsi|supér.}), et
